% problem_id: S001-lemma-cartesian-colimit-kappa
% source_id: S001 (Stacks Project)
% source_file: sites-cohomology.tex
% statement_locator: sites-cohomology.tex lines 11295--11309
% proof_locator: sites-cohomology.tex lines 11311--11455
% env: lemma
% id_note: label 在本来源内唯一
% pairing: tight (verified)
% status: complete (statement + author proof verbatim + reason + locators)
%
\begin{lemma}
\label{lemma-cartesian-colimit-kappa}
In the situation above, there exists a cardinal $\kappa$ with the following
properties:
\begin{enumerate}
\item for every nonzero object $K$ of $\mathit{QC}(\mathcal{O})$ there exists a
nonzero morphism $E \to K$ of $\mathit{QC}(\mathcal{O})$ such that
$|E| \leq \kappa$,
\item for every morphism $\alpha : E \to \bigoplus_n K_n$ of
$\mathit{QC}(\mathcal{O})$ such that $|E| \leq \kappa$, there
exist morphisms $E_n \to K_n$
in $\mathit{QC}(\mathcal{O})$ with $|E_n| \leq \kappa$
such that $\alpha$ factors through $\bigoplus E_n \to \bigoplus K_n$.
\end{enumerate}
\end{lemma}

\begin{proof}
Let $\kappa$ be an upper bound for the following set of cardinals:
\begin{enumerate}
\item $|\coprod_V j_{U!}\mathcal{O}_U(V)|$ for all $U \in \Ob(\mathcal{C})$,
\item the cardinals $\kappa(\mathcal{O}(V) \to \mathcal{O}(U))$
found in More on Algebra, Lemma \ref{more-algebra-lemma-enlarge}
for all morphisms $U \to V$ in $\mathcal{C}$,
\item the cardinal found in Lemma \ref{lemma-qis-small}.
\end{enumerate}
We claim that for any complex $\mathcal{F}^\bullet$ representing an object of
$\mathit{QC}(\mathcal{O})$ and any subcomplex
$\mathcal{S}^\bullet \subset \mathcal{F}^\bullet$
with $|\mathcal{S}^\bullet| \leq \kappa$ there exists a subcomplex
$\mathcal{H}^\bullet$ of $\mathcal{F}^\bullet$ containing $\mathcal{S}^\bullet$
such that $\mathcal{H}^\bullet$ represents an object of
$\mathit{QC}(\mathcal{O})$
and such that $|\mathcal{H}^\bullet| \leq \kappa$.
In the next two paragraphs we show that the claim implies the lemma.

\medskip\noindent
As in (1) let $K$ be a nonzero object of $\mathit{QC}(\mathcal{O})$.
Say $K$ is represented by the complex
of $\mathcal{O}$-modules $\mathcal{F}^\bullet$. Then $H^i(\mathcal{F}^\bullet)$
is nonzero for some $i$. Hence there exists an object $U$ of $\mathcal{C}$
and a section $s \in \mathcal{F}^i(U)$ with $d(s) = 0$ which determines a
nonzero section of $H^i(\mathcal{F}^\bullet)$ over $U$. Then the image
of $s : j_{U!}\mathcal{O}_U[-i] \to \mathcal{F}^\bullet$ is a subcomplex
$\mathcal{S}^\bullet \subset \mathcal{F}^\bullet$ with
$|\mathcal{S}^\bullet| \leq \kappa$. Applying the claim we get
$\mathcal{H}^\bullet \to \mathcal{F}^\bullet$ in
$\mathit{QC}(\mathcal{O})$ nonzero
with $|\mathcal{H}^\bullet| \leq \kappa$. Thus (1) holds.

\medskip\noindent
Let $\alpha : E \to \bigoplus K_n$ be as in (2).
Choose any complexes $\mathcal{K}_n^\bullet$ representing $K_n$.
Then $\bigoplus \mathcal{K}_n^\bullet$ represents $\bigoplus K_n$.
By the construction of the derived category we can represent
$E$ by a complex $\mathcal{E}^\bullet$ such that $\alpha$
is represented by a morphism
$a : \mathcal{E}^\bullet \to \bigoplus \mathcal{K}_n^\bullet$
of complexes. By Lemma \ref{lemma-qis-small} and our choice of
$\kappa$ above we may assume
$|\mathcal{E}^\bullet| \leq \kappa$. By the claim
we get subcomplexes $\mathcal{E}_n^\bullet \subset \mathcal{K}_n^\bullet$
representing objects $E_n$ of
$\mathit{QC}(\mathcal{O})$ with $|E_n| \leq \kappa$
containing the image of $a_n : \mathcal{E}^\bullet \to \mathcal{K}_n^\bullet$
as desired.

\medskip\noindent
Proof of the claim. Let $\mathcal{F}^\bullet$ be a complex representing
an object of $\mathit{QC}(\mathcal{O})$ and let
$\mathcal{S}^\bullet \subset \mathcal{F}^\bullet$ be a subcomplex of
size $\leq \kappa$. We are going to inductively construct subcomplexes
$$
\mathcal{S}^\bullet = \mathcal{S}_0^\bullet \subset
\mathcal{S}_1^\bullet \subset
\mathcal{S}_2^\bullet \subset \ldots \subset \mathcal{F}^\bullet
$$
of size $\leq \kappa$ such that for every morphism $f : U \to V$
of $\mathcal{C}$ and every $i \in \mathbf{Z}$
\begin{enumerate}
\item the kernel of the arrow
$H^i(\mathcal{S}_n^\bullet(V)
\otimes_{\mathcal{O}(V)}^\mathbf{L} \mathcal{O}(U)) \to
H^i(\mathcal{S}_n^\bullet(U))$
maps to zero in
$H^i(\mathcal{S}_{n + 1}^\bullet(V)
\otimes_{\mathcal{O}(V)}^\mathbf{L} \mathcal{O}(U))$,
\item the image of the arrow
$H^i(\mathcal{S}_n^\bullet(U)) \to H^i(\mathcal{S}_{n + 1}^\bullet(U))$
is contained in the image of
$H^i(\mathcal{S}_{n + 1}^\bullet(V)
\otimes_{\mathcal{O}(V)}^\mathbf{L} \mathcal{O}(U)) \to
H^i(\mathcal{S}_{n + 1}^\bullet(U))$,
\end{enumerate}
Once this is done we can set
$\mathcal{H}^\bullet = \bigcup \mathcal{S}_n^\bullet$.
Namely, since derived tensor product and taking cohomology of
complexes of modules over rings commute with filtered
colimits, the conditions (1) and (2) together will guarantee
that
$$
\mathcal{H}^\bullet(V) \otimes_{\mathcal{O}(V)}^\mathbf{L} \mathcal{O}(U)
\longrightarrow
\mathcal{H}^\bullet(U)
$$
is an isomorphism on cohomology in all degrees and hence an isomorphism
in $D(\mathcal{O}(U))$ for all $f : U \to V$ in $\mathcal{C}$.
Hence $\mathcal{H}^\bullet$ represents an object of
$\mathit{QC}(\mathcal{O})$ as desired.

\medskip\noindent
Construction of $\mathcal{S}_{n + 1}$ given $\mathcal{S}_n$.
For every morphism $f : U \to V$ of $\mathcal{C}$
we consider the commutative diagram
$$
\xymatrix{
\mathcal{S}_n^\bullet(V)
\ar[r] \ar[d] &
\mathcal{S}_n^\bullet(U) \ar[d] \\
\mathcal{F}^\bullet(V)
\ar[r] &
\mathcal{F}^\bullet(U)
}
$$
This is a diagram as in More on Algebra, Lemma \ref{more-algebra-lemma-enlarge}
for the ring map $\mathcal{O}(V) \to \mathcal{O}(U)$, i.e., the bottom row
induces an isomorphism
$$
\mathcal{F}^\bullet(V) \otimes_{\mathcal{O}(V)}^\mathbf{L} \mathcal{O}(U)
\longrightarrow
\mathcal{F}^\bullet(U)
$$
in $D(\mathcal{O}(U))$. Thus we may choose subcomplexes
$$
\mathcal{S}_n^\bullet(V) \subset M^\bullet_f \subset \mathcal{F}^\bullet(V)
\quad\text{and}\quad
\mathcal{S}_n^\bullet(U) \subset N^\bullet_f \subset \mathcal{F}^\bullet(U)
$$
as in More on Algebra, Lemma \ref{more-algebra-lemma-enlarge}
and in particular we see that $|N^i_f|, |M^i_f| \leq \kappa$. Next, we apply
Lemma \ref{lemma-build-from} using the subsets
$$
\mathcal{S}_n^i(U) \amalg \coprod\nolimits_{f : U \to V} N^i_f
\amalg \coprod\nolimits_{g : W \to U} M^i_g
\subset
\mathcal{F}^i(U)
$$
to find a subcomplex
$$
\mathcal{S}_n^\bullet \subset
\mathcal{S}_{n + 1}^\bullet \subset
\mathcal{F}^\bullet
$$
with containing those subsets and such that
$|\mathcal{S}_{n + 1}^\bullet| \leq \kappa$.
Conditions (1) and (2) hold because the corresponding statements
hold for 
$\mathcal{S}_n^\bullet(V) \subset M^\bullet_f$ and
$\mathcal{S}_n^\bullet(U) \subset N^\bullet_f$ by the construction in
More on Algebra, Lemma \ref{more-algebra-lemma-enlarge}.
Thus the proof is complete.
\end{proof}
